\scanpage{139}
\margin{\deleted{(où $D'^{*}$ est l'image inv. de $\Theta^{(d_1)*}$)}}
fondamentaux $\Pi_1P_{D''^{*}},\Pi_1P_{D'^{*}}$, Cette fois, on a un morphisme canonique (dit ``de générisation'')
\[
(231)\qquad \boxed{\Pi_1P_{D''^{*}}\xrightarrow{\gamma_{D'',D'}}\Pi_1P_{D'^{*}}}.
\]
\deleted{\unsure{compatible avec les opérations de}}

On donnera la \inserted{\unsure{principe de|principale}} construction de ce morphisme, et ses principales propriétés formelles, dans la suite. Notons tout de suite qu'on a un morphisme trivial, déduit de l'inclusion
\[
\deleted{\unsure{P_{D'^{*}}=P_{D'}\vert D'^{*}\longrightarrow P_{D'}}}
\]
\[
P_{D'^{*}}=P_{D'}\vert D'^{*}\lhook\joinrel\longrightarrow P_{D'}
\]
\[
(232)\qquad\Pi_1P_{D'^{*}}\longrightarrow\Pi_1P_{D'}
\]
donc le composé de (231) et (232) fournit
\[
(233)\qquad\Pi_1P_{D'^{*}}\longrightarrow\Pi_1P_{D'},
\]
qui, lui, est de nature essentiellement trivial, car on a en effet un morphisme évident $P_{D'}\longrightarrow P_{D''}$. En fait,
\margin{138}
